\begin{prob}
	Suppose we are given both an undirected graph G with weighted edges and a minimum spanning tree T of G.
	\begin{subprobset}
	    \begin{subprob}
    		Describe an efficient algorithm to update the minimum spanning tree when the weight of one edge e in T is decreased.
    		\begin{responsebox}{1in}
    		    The minimum spanning tree of the updated graph would be T.
    		\end{responsebox}
	    \end{subprob}

        \newpage

	    \begin{subprob}
    		Describe an efficient algorithm to update the minimum spanning tree when the weight of one edge e not in T is increased.
    		\begin{responsebox}{1in}
    		    The minimum spanning tree of the updated graph would be T.
    		\end{responsebox}
	    \end{subprob}

  	    \begin{subprob}
    		Describe an efficient algorithm to update the minimum spanning tree when the weight of one edge e in T is increased.
    		\begin{responsebox}{2in}
    		    Let $e = (u,v)$ be the edge whose weight is increased. Remove the edge e from the minimum spanning tree T.
    		    This divides the tree T into two connected components. Let $T_u$ be the component which contains u and $T_v$ be the component which contains v.
    		    We can identify $T_u$ and $T_v$ by running a BFS or DFS with u and v as the sources. This takes time $\Theta(V + E)$.
    		    While running BFS we can also label each node as 0 if it is a part of $T_u$ and 1 if it is a part of $T_v$.
    		    We can now examine each edge e in the graph and find the minimum weight edge which connects a node labelled 0 to a node labelled 1.
    		    This takes time $\Theta(E)$.
    		    We can then add this edge to T to get the minimum spanning tree of the updated graph. The total time complexity is $\Theta(V + E)$.
    		\end{responsebox}
	    \end{subprob}


	    \begin{subprob}
    		Describe an efficient algorithm to update the minimum spanning tree when the weight of one edge e not in T is decreased.
    		\begin{responsebox}{2in}
    		    Let $e = (u,v)$ be the edge whose weight is decreased. Add the edge e to the minimum spanning tree T. This would result in a cycle in T.
    		    We can identify the nodes and the edges on the cycle by running BFS or DFS, with minimal modifications, with u or v as the source. This takes time $\Theta(V + E)$.
    	            Find the maximum weight edge on the cycle and remove it from T. This takes time $\Theta(E)$.
    		    The resultant tree would be a minimum spanning tree for the updated graph. The total time complexity is $\Theta(V + E)$.
    		\end{responsebox}
        \end{subprob}
    \end{subprobset}
\end{prob}